\documentclass[12pt,a4paper]{article}

\usepackage[a4paper,margin=1in]{geometry}
\usepackage{graphicx}
\usepackage{booktabs}
\usepackage{amsmath,amssymb}
\usepackage{float}
\usepackage{hyperref}
\usepackage{xcolor}
\usepackage{listings}
\usepackage{enumitem}
\usepackage{fancyhdr}
\usepackage{longtable}

\hypersetup{colorlinks=true,linkcolor=blue,urlcolor=blue}

\pagestyle{fancy}
\fancyhf{}
\lhead{ICS1512}
\rhead{Machine Learning Algorithms Laboratory}
\cfoot{\thepage}

\lstset{
language=Python,
basicstyle=\ttfamily\small,
numbers=left,
frame=single,
breaklines=true,
showstringspaces=false,
keywordstyle=\color{blue},
commentstyle=\color{gray},
stringstyle=\color{purple}
}

\begin{document}

\begin{tabular}{|p{4cm}|p{11cm}|}
\hline
Experiment No. & 2 \\ \hline
Experiment Title & Email Spam/Ham Classification using Na\"ive Bayes and KNN
\footnote{\textbf{GitHub Repository: }\url{https://github.com/Nightingales21/ICS1512-Machine-Learning-Labratory}}\\ \hline
Student Name & \\ \hline
Register Number & \\ \hline
Faculty & Dr. Poreddy Ajay Kumar Reddy \\ \hline
Submission Date & \\ \hline
\end{tabular}

\section{Objective}
To build and compare email spam/ham classifiers using Na\"ive Bayes and
$k$-Nearest Neighbors (KNN). Specifically: (i) compare Gaussian, Multinomial,
and Bernoulli Na\"ive Bayes variants; (ii) study the effect of varying $k$ on
KNN performance; (iii) compare the KDTree and BallTree neighbor-search
algorithms; (iv) optimize KNN hyperparameters using \texttt{GridSearchCV} and
\texttt{RandomizedSearchCV}; (v) evaluate theoretical versus experimental
computational complexity and execution time; and (vi) validate results using
5-fold cross-validation.

\section{Problem Statement}
Given the Spambase dataset of 4{,}601 emails, each represented by 57
numeric features (word/character frequencies and capital-letter run-length
statistics) and a binary label (1 = spam, 0 = ham), the goal is to train
classifiers that correctly predict whether a given email is spam. This is a
binary supervised classification problem. Both a generative probabilistic
approach (Na\"ive Bayes, in three variants) and a non-parametric,
distance-based approach (KNN) are built, tuned, and compared on accuracy,
precision, recall, F1-score, ROC-AUC, and computational efficiency.

\section{Dataset Description}
\begin{longtable}{|p{6cm}|p{8cm}|}
\hline
Dataset Name & Spambase \\ \hline
Dataset Source & UCI Machine Learning Repository / Kaggle
(\url{https://www.kaggle.com/datasets/somesh24/spambase}) \\ \hline
Number of Samples & 4{,}601 \\ \hline
Number of Features & 57 (48 word-frequency, 6 character-frequency, 3
capital-run-length features) \\ \hline
Number of Classes & 2 (Spam = 1, Ham = 0) \\ \hline
Missing Values & 0 (verified programmatically) \\ \hline
Train--Test Split & 80\% / 20\%, stratified (additional splits 90/10, 70/30,
60/40 evaluated in Section~\ref{sec:additional}) \\ \hline
\end{longtable}

Class balance: 2{,}788 ham emails (60.6\%) and 1{,}813 spam emails (39.4\%) --
a moderate, not severe, class imbalance.

\section{Exploratory Data Analysis}
Include all relevant visualizations.
\begin{itemize}
\item Dataset overview
\item Statistical summary
\item Missing value analysis
\item Class distribution
\item Correlation matrix
\item Feature distribution
\end{itemize}

The reusable function \texttt{generate\_eda\_summary()} from
\texttt{ml\_lab\_utils.py} (developed in Experiment 1) was used unmodified to
produce the consolidated 12-panel EDA figure below.

\begin{figure}[H]
\centering
\includegraphics[width=\linewidth]{figures/eda_spambase.png}
\caption{Consolidated 12-panel EDA summary for the Spambase dataset.}
\label{fig:eda}
\end{figure}

\textbf{Inference:} The dataset has no missing values but contains 391
duplicate rows, consistent with repeated boilerplate/template spam emails.
The class distribution confirms a moderate imbalance (60.6\% ham, 39.4\%
spam). The two highest-variance numeric features are
\texttt{capital\_run\_length\_total} and \texttt{capital\_run\_length\_longest},
both extremely right-skewed (visible in Panels 6, 9, 10, and 12) --- long,
heavy-tailed runs of capital letters are a strong, rare-but-informative spam
signal. The Q--Q plot (Panel 10) confirms these features are far from
normally distributed, which motivates feature scaling before use with
distance-based (KNN) and Gaussian-likelihood (Gaussian NB) models. Most
individual word/character-frequency features are dominated by zeros (sparse),
which favors Bernoulli and Multinomial Na\"ive Bayes formulations that are
designed for exactly this kind of frequency/occurrence data.

\section{Data Preprocessing}

\begin{itemize}
\item \textbf{Missing value handling:} None required; the dataset was
verified to contain zero missing cells.
\item \textbf{Encoding:} Not required; all 57 features and the target are
already numeric.
\item \textbf{Feature scaling:} \texttt{StandardScaler} (fit on the training
split only) was applied before Gaussian NB and all KNN variants, since both
are sensitive to feature scale. Multinomial NB and Bernoulli NB were trained
on the original (unscaled, non-negative) frequency features, since
Multinomial NB requires non-negative counts/frequencies and Bernoulli NB
internally binarizes at a threshold.
\item \textbf{Train--test split:} \texttt{train\_test\_split} with
\texttt{test\_size=0.2}, \texttt{stratify=y}, \texttt{random\_state=42}.
\item \textbf{Feature engineering:} Not required; the 57 pre-computed
frequency/run-length features from the original Spambase specification were
used directly.
\end{itemize}

\section{Machine Learning Algorithm}

\textbf{Na\"ive Bayes.} A probabilistic classifier based on Bayes' theorem
with the ``na\"ive'' assumption of conditional independence between features
given the class:
\[
P(C \mid X) = \frac{P(X \mid C)\,P(C)}{P(X)}
\]
\begin{itemize}
\item \textit{Gaussian NB} assumes each feature is normally distributed
within each class.
\item \textit{Multinomial NB} models feature vectors as multinomial counts
(natural fit for word-frequency data).
\item \textit{Bernoulli NB} models each feature as a binary indicator
(presence/absence), which is well suited to sparse, mostly-zero
word/character-frequency vectors.
\end{itemize}
\textbf{Advantages:} Fast to train (single pass), works well with
high-dimensional/sparse data, requires little training data.
\textbf{Limitations:} The conditional-independence assumption is rarely
exactly true; Gaussian NB is sensitive to strongly non-normal features (as
seen in Section 3).

\textbf{K-Nearest Neighbors (KNN).} A lazy, instance-based learning
algorithm that classifies a query point by majority vote among its $k$
nearest neighbors in feature space, using a distance metric (Euclidean or
Manhattan) and, optionally, distance-weighted voting.
\textbf{Advantages:} Simple, non-parametric, no explicit training phase,
naturally handles multi-class problems.
\textbf{Limitations:} Prediction is computationally expensive at scale
(especially with brute-force search), sensitive to feature scaling and the
curse of dimensionality, and requires careful tuning of $k$, the distance
metric, and the neighbor-search algorithm.

\section{Experimental Setup}
\begin{tabular}{ll}
\toprule
Python Version & 3.12 \\
Scikit-Learn Version & 1.8.0 \\
IDE & Jupyter Notebook / VS Code \\
Operating System & Ubuntu 24.04 LTS \\
Hardware & Standard lab workstation (CPU only) \\
\bottomrule
\end{tabular}

\vspace{0.3cm}
\textbf{Environment activation command (as mandated by the lab manual):}
\begin{lstlisting}[language=bash]
source /opt/anaconda3/bin/activate anaconda-navigation
\end{lstlisting}

\section{Hyperparameter Selection}

\begin{tabular}{ll}
\toprule
Parameter & Value\\
\midrule
\texttt{test\_size} & 0.2 \\
\texttt{random\_state} & 42 \\
KNN \texttt{n\_neighbors} search space & \{1, 3, 5, 7, 9, 11, 13, 15\} \\
KNN \texttt{weights} search space & \{uniform, distance\} \\
KNN \texttt{metric} search space & \{euclidean, manhattan\} \\
KNN \texttt{algorithm} search space & \{kd\_tree, ball\_tree\} \\
Cross-validation folds (tuning) & 5 (Stratified K-Fold) \\
RandomizedSearchCV \texttt{n\_iter} & 20 \\
\bottomrule
\end{tabular}

\section{Code Organization}
Per the lab manual (Section 4), all reusable functionality lives in
\texttt{ml\_lab\_utils.py} (unchanged from Experiment 1) and is imported
directly:
\begin{itemize}
\item \texttt{generate\_eda\_summary()} -- 12-panel EDA figure (Section 3).
\item \texttt{train\_evaluate\_classification()} / \\
\texttt{classification\_performance\_metrics()} -- classifier
training/evaluation and metric computation, reused for every model in this
experiment.
\item \texttt{set\_plot\_style()} -- global Times New Roman / 15pt / bold-axis
formatting applied to every figure below.
\end{itemize}

\begin{lstlisting}[language=Python, caption={Naive Bayes variants -- driver excerpt from \texttt{experiment2\_spambase.py}.}]
gnb = GaussianNB()
gnb.fit(X_train_scaled, y_train)
y_pred = gnb.predict(X_test_scaled)

mnb = MultinomialNB()
mnb.fit(X_train_raw, y_train)   # requires non-negative features

bnb = BernoulliNB()
bnb.fit(X_train_raw, y_train)   # internally binarizes features
\end{lstlisting}

\begin{lstlisting}[language=Python, caption={KNN hyperparameter tuning excerpt.}]
param_grid = {
    "n_neighbors": [1, 3, 5, 7, 9, 11, 13, 15],
    "weights": ["uniform", "distance"],
    "metric": ["euclidean", "manhattan"],
    "algorithm": ["kd_tree", "ball_tree"],
}
grid_search = GridSearchCV(KNeighborsClassifier(), param_grid,
                            cv=StratifiedKFold(5, shuffle=True, random_state=42),
                            scoring="accuracy", n_jobs=-1)
grid_search.fit(X_train_scaled, y_train)
\end{lstlisting}

The full source of \texttt{ml\_lab\_utils.py} and
\texttt{experiment2\_spambase.py} is included in the accompanying GitHub
repository (see title-page footnote) and in the submitted code archive.

\section{Results}

\textbf{Na\"ive Bayes Comparison}
\begin{tabular}{lccccc}
\toprule
Metric & Gaussian & Multinomial & Bernoulli \\
\midrule
Accuracy  & 0.8328 & 0.7763 & \textbf{0.8762} \\
Precision & 0.8666 & 0.7757 & 0.8760 \\
Recall    & 0.8328 & 0.7763 & 0.8762 \\
F1        & 0.8345 & 0.7760 & 0.8753 \\
ROC-AUC   & 0.9376 & 0.8248 & \textbf{0.9496} \\
\bottomrule
\end{tabular}

\vspace{0.3cm}
\textbf{KNN Comparison (varying $k$, Euclidean, uniform weights, scaled features)}
\begin{tabular}{lcccc}
\toprule
$k$ & Accuracy & Precision & Recall & F1 \\
\midrule
1  & 0.8990 & 0.8792 & 0.8623 & 0.8707 \\
3  & 0.8979 & 0.8726 & 0.8678 & 0.8702 \\
5  & 0.9077 & 0.8861 & 0.8788 & 0.8824 \\
7  & 0.9088 & 0.8930 & 0.8733 & 0.8830 \\
9  & 0.9088 & 0.8930 & 0.8733 & 0.8830 \\
11 & \textbf{0.9099} & \textbf{0.9000} & 0.8678 & \textbf{0.8836} \\
\bottomrule
\end{tabular}

\vspace{0.3cm}
\textbf{GridSearchCV vs. RandomizedSearchCV}
\begin{tabular}{lcc}
\toprule
Parameter & GridSearchCV & RandomizedSearchCV \\
\midrule
Best $k$        & 9         & 9 \\
Metric          & manhattan & manhattan \\
Weights         & distance  & distance \\
Algorithm       & kd\_tree  & kd\_tree \\
CV Accuracy     & 0.9253    & 0.9253 \\
Execution Time  & 30.52 s   & 9.36 s \\
\bottomrule
\end{tabular}

\vspace{0.3cm}
\textbf{KDTree vs. BallTree} (best tuned KNN: $k=9$, manhattan, distance-weighted)
\begin{tabular}{lcc}
\toprule
Metric & KDTree & BallTree \\
\midrule
Accuracy         & 0.9207 & 0.9207 \\
Training Time    & 0.0136 s & 0.0079 s \\
Prediction Time  & 0.1622 s & 0.1252 s \\
\bottomrule
\end{tabular}

\vspace{0.3cm}
\textbf{Cross Validation (5-fold, Bernoulli NB vs. Best Tuned KNN)}
\begin{tabular}{lcc}
\toprule
Fold & Na\"ive Bayes (Bernoulli) & Best KNN \\
\midrule
1 & 0.8849 & 0.9164 \\
2 & 0.8859 & 0.9326 \\
3 & 0.8793 & 0.9359 \\
4 & 0.9022 & 0.9217 \\
5 & 0.8815 & 0.9217 \\
\midrule
Average & 0.8868 & \textbf{0.9257} \\
\bottomrule
\end{tabular}

\vspace{0.3cm}
\textbf{Theoretical Time Complexity}
\begin{tabular}{lcc}
\toprule
Algorithm & Training & Prediction \\
\midrule
Na\"ive Bayes  & $O(nd)$      & $O(d)$ \\
KNN (Brute)   & $O(1)$       & $O(nd)$ \\
KDTree        & $O(n\log n)$ & $O(\log n)_{avg}$ \\
BallTree      & $O(n\log n)$ & $O(\log n)_{avg}$ \\
\bottomrule
\end{tabular}

\vspace{0.3cm}
\textbf{Experimental Time Analysis}
\begin{tabular}{lcc}
\toprule
Algorithm & Training (s) & Prediction (s) \\
\midrule
Gaussian NB     & 0.0118 & 0.00070 \\
Multinomial NB  & 0.0020 & 0.00023 \\
Bernoulli NB    & 0.0130 & 0.00094 \\
Best KNN (tuned)& 0.0136 (kd\_tree) & 0.1622 (kd\_tree) \\
\bottomrule
\end{tabular}

\vspace{0.3cm}
\textbf{Overall Best Model (Test Set): Best Tuned KNN} --- Accuracy 0.9207,
Precision 0.9220, Recall 0.9207, F1 0.9199, ROC-AUC 0.9712.

\section{Result Visualization}
\label{sec:additional}

For \textbf{every graph}, a figure, caption, and inference are provided.

\begin{figure}[H]
\centering
\includegraphics[width=0.75\linewidth]{figures/accuracy_vs_k.png}
\caption{Accuracy vs. $k$ for plain (untuned) KNN.}
\label{fig:acc_k}
\end{figure}
\textbf{Inference:} Accuracy rises from $k=1$ to $k=11$, with the sharpest
jump between $k=3$ and $k=5$. Very small $k$ (particularly $k=1$) is more
sensitive to noisy individual training points (overfitting), while the
curve's plateau beyond $k=7$ suggests diminishing returns and the onset of
underfitting for excessively large $k$.

\begin{figure}[H]
\centering
\includegraphics[width=0.75\linewidth]{figures/gridsearch_heatmap.png}
\caption{GridSearchCV mean CV accuracy across $k$ and distance metric.}
\label{fig:grid_heatmap}
\end{figure}
\textbf{Inference:} Manhattan distance consistently outperforms Euclidean
distance across nearly all values of $k$ in cross-validation, which is
consistent with the sparse, frequency-based nature of the Spambase features
(Manhattan distance is generally more robust to high-dimensional sparse
data than Euclidean distance).

\begin{figure}[H]
\centering
\includegraphics[width=0.75\linewidth]{figures/randomsearch_score_distribution.png}
\caption{Distribution of mean CV accuracy scores sampled by RandomizedSearchCV.}
\label{fig:random_dist}
\end{figure}
\textbf{Inference:} Most of the 20 randomly sampled hyperparameter
combinations cluster around 0.88--0.92 accuracy, with a small number of
combinations reaching the same optimum (0.9253) found by the exhaustive grid
search -- confirming that random search found the global optimum in this
search space using roughly 3$\times$ fewer fit calls.

\begin{figure}[H]
\centering
\includegraphics[width=\linewidth]{figures/roc_curves.png}
\caption{ROC curves for Gaussian NB, Multinomial NB, Bernoulli NB, and the best tuned KNN.}
\label{fig:roc}
\end{figure}
\textbf{Inference:} The tuned KNN model dominates all three Na\"ive Bayes
variants across the full ROC curve (AUC = 0.971), followed closely by
Bernoulli NB (AUC = 0.950). Multinomial NB has the weakest separability
(AUC = 0.825), consistent with its lower accuracy and its independence
assumption being a poorer fit for continuous frequency values compared to
Bernoulli's binary presence/absence framing.

\begin{figure}[H]
\centering
\includegraphics[width=\linewidth]{figures/precision_recall_curves.png}
\caption{Precision--Recall curves for all four classifiers.}
\label{fig:pr}
\end{figure}
\textbf{Inference:} The best tuned KNN maintains the highest precision across
almost the entire recall range, meaning it is the most reliable model when
minimizing false-positive (ham misclassified as spam) errors is a priority
-- an important practical consideration for a real spam filter.

\begin{figure}[H]
\centering
\includegraphics[width=\linewidth]{figures/cm_best_knn.png}
\caption{Confusion matrix for the best tuned KNN model on the held-out test set.}
\label{fig:cm_knn}
\end{figure}
\textbf{Inference:} The confusion matrix confirms strong performance on both
classes with balanced error rates; most confusion occurs on borderline
emails whose word/character frequencies fall between typical spam and ham
patterns, which is also reflected in the KDE overlap visible in Panel 11 of
the EDA figure.

\begin{figure}[H]
\centering
\includegraphics[width=\linewidth]{figures/classifier_comparison_bar.png}
\caption{Overall test-set accuracy comparison across all four classifiers.}
\label{fig:bar}
\end{figure}
\textbf{Inference:} The tuned KNN model is the clear winner on accuracy,
followed by Bernoulli NB, Gaussian NB, and finally Multinomial NB. This
ordering matches the ROC-AUC ranking exactly, giving consistent evidence
across two independent metrics for the same relative model ranking.

\begin{figure}[H]
\centering
\includegraphics[width=\linewidth]{figures/cv_accuracy.png}
\caption{5-fold cross-validation accuracy: Bernoulli NB vs. best tuned KNN.}
\label{fig:cv}
\end{figure}
\textbf{Inference:} KNN's cross-validation accuracy is both higher on average
(0.9257 vs. 0.8868) and more stable across folds than Bernoulli NB, giving
confidence that KNN's advantage on the single train/test split is not an
artifact of that particular split.

\begin{figure}[H]
\centering
\includegraphics[width=\linewidth]{figures/training_prediction_time.png}
\caption{Training and prediction time comparison across the three Na\"ive Bayes variants.}
\label{fig:time}
\end{figure}
\textbf{Inference:} Multinomial NB is fastest to both train and predict,
owing to its simple count-based sufficient statistics; Bernoulli and
Gaussian NB are slightly slower due to the additional per-feature
binarization/likelihood computations, but all three remain on the order of
milliseconds, confirming Na\"ive Bayes's practical suitability for
large-scale, low-latency spam filtering.

\subsection*{Additional Tasks}

\textbf{Different train-test splits (Best tuned KNN):}
\begin{tabular}{lc}
\toprule
Test Size & Accuracy \\
\midrule
0.10 & 0.9067 \\
0.20 & 0.9207 \\
0.30 & \textbf{0.9312} \\
0.40 & 0.9202 \\
\bottomrule
\end{tabular}
\quad Accuracy is fairly stable (0.907--0.931) across all four splits,
indicating the model generalizes well and is not overly sensitive to the
exact train/test proportion.

\vspace{0.3cm}
\textbf{Euclidean vs. Manhattan distance ($k=9$):}
\begin{tabular}{lc}
\toprule
Distance & Accuracy \\
\midrule
Euclidean & 0.9088 \\
Manhattan & 0.9034 \\
\bottomrule
\end{tabular}
\quad On the single 80/20 test split Euclidean is marginally ahead, but the
5-fold GridSearchCV result (which averages over many splits) favored
Manhattan -- the cross-validated result is the more reliable indicator of
true generalization performance.

\vspace{0.3cm}
\textbf{Weighted KNN ($k=9$):}
\begin{tabular}{lc}
\toprule
Weighting & Accuracy \\
\midrule
Uniform  & 0.9088 \\
Distance & \textbf{0.9164} \\
\bottomrule
\end{tabular}
\quad Distance-weighting improves accuracy by giving closer neighbors more
influence over the vote, reducing the impact of borderline/noisy far
neighbors -- consistent with distance-weighting being selected by both
GridSearchCV and RandomizedSearchCV as part of the overall best
configuration.

\vspace{0.3cm}
\textbf{Overfitting/underfitting discussion:} Small $k$ (e.g., $k=1$) fits
the training data almost perfectly but is highly sensitive to individual
noisy points, leading to higher variance and mild overfitting -- visible as
lower and less stable accuracy at $k=1$ and $k=3$ in
Figure~\ref{fig:acc_k}. Large $k$ averages over many neighbors, smoothing
the decision boundary and increasing bias; the plateauing/slight-decline
pattern typically observed for very large $k$ (beyond the tested range here)
would indicate underfitting, as the neighborhood becomes large enough to mix
in points from the opposite class.

\section{Discussion}
Across every evaluation lens used in this experiment -- single-split
accuracy/precision/recall/F1, ROC-AUC, 5-fold cross-validation, and
precision-recall curves -- the tuned KNN classifier consistently
outperformed all three Na\"ive Bayes variants. Among the Na\"ive Bayes
family, Bernoulli NB was clearly the strongest, which is expected given the
sparse, largely binary-informative nature of word/character-presence
features in spam detection; Multinomial NB underperformed relative to
Bernoulli NB here, and Gaussian NB was hurt by the strong non-normality of
several features (confirmed in the EDA Q--Q plot). GridSearchCV and
RandomizedSearchCV converged on the identical optimal hyperparameter
combination (manhattan distance, distance-weighted voting, $k=9$),
demonstrating that RandomizedSearchCV can match exhaustive search quality at
roughly 3.3$\times$ lower computational cost on this search space. KDTree
and BallTree produced identical accuracy (as expected, since both are exact
nearest-neighbor algorithms), with BallTree marginally faster in both
training and prediction on this dataset's scale. The theoretical complexity
table is consistent with the experimental timings: Na\"ive Bayes's
$O(nd)$/$O(d)$ training/prediction cost is essentially instantaneous, while
KNN's tree-based prediction, though sub-linear in theory, carries a larger
constant-factor overhead than Na\"ive Bayes at this dataset size (4{,}601
samples), which is why the theoretical asymptotic advantage of tree-based
KNN prediction does not yet dominate in wall-clock time on this data.

\section{Conclusion}
This experiment successfully built, tuned, and compared Na\"ive Bayes
(Gaussian, Multinomial, Bernoulli) and KNN classifiers for email spam
detection on the Spambase dataset. Bernoulli Na\"ive Bayes was the best
Na\"ive Bayes variant (Accuracy 0.876, ROC-AUC 0.950), consistent with the
dataset's sparse, presence/absence-style features. The optimal KNN
configuration, identified identically by GridSearchCV and
RandomizedSearchCV, used $k=9$, Manhattan distance, and distance-weighted
voting, achieving the best overall test performance (Accuracy 0.921, ROC-AUC
0.971) and the best 5-fold cross-validation accuracy (0.926). KDTree and
BallTree were equivalent in accuracy with BallTree marginally faster.
For large-scale or latency-sensitive spam filtering, Na\"ive Bayes remains
attractive due to its near-constant, sub-millisecond training and prediction
time; when maximum accuracy is the priority and moderate prediction latency
is acceptable, the tuned KNN classifier is preferred based on the results of
this experiment.

\section{References}
\begin{enumerate}
\item M. Hopkins, E. Reeber, G. Forman, and J. Suermondt, ``Spambase Data
Set,'' UCI Machine Learning Repository, 1999. [Online]. Available:
\url{https://archive.ics.uci.edu/dataset/94/spambase}
\item F. Pedregosa et al., ``Scikit-learn: Machine Learning in Python,''
\emph{Journal of Machine Learning Research}, vol. 12, pp. 2825--2830, 2011.
\item C. M. Bishop, \emph{Pattern Recognition and Machine Learning}. Springer,
2006.
\item A. G\'eron, \emph{Hands-On Machine Learning with Scikit-Learn, Keras
and TensorFlow}, 3rd ed. O'Reilly Media, 2022.
\end{enumerate}

\end{document}
